\documentclass[11pt]{letter}

\usepackage[margin=0.5in]{geometry}
\usepackage{hyperref}

\begin{document}

\begin{letter}{}

\opening{Dear PLATO Hiring Team,}

I am writing to apply for the Junior Data Engineer position with PLATO, Job Number RR\#007. With a Master of Science in Computer Science from the University of Calgary and hands-on experience building Python data-processing systems, backend APIs, ETL-style workflows, and scalable research software, I am excited by the opportunity to contribute to reliable data solutions for PLATO's client projects.

What stands out to me about PLATO is the combination of technical delivery, quality, mentorship, and inclusive impact. I like that this role is junior-friendly while still connected to real data needs: creating pipelines, supporting ETL processes, improving accessibility and usability of data, troubleshooting reliability issues, and working with cross-functional teams. That mix fits how I like to work: practical, careful, and focused on making data systems understandable and dependable.

As a Research Software Developer with BigGeo/Mitacs at the University of Calgary, I built C++ and Python systems for processing, storing, and retrieving large Earth observation datasets. I developed repeatable pipelines, automated satellite-imagery preprocessing, converted data into structured tensor files, and improved one encoding workflow from 233 seconds to 26 seconds using C++ multithreading. This work strengthened my ability to reason about data structures, scalability, data quality, performance, and reliability in high-volume data systems.

I also bring backend and database experience that connects well with data engineering. As a Back-End Developer, I designed and implemented Django REST APIs backed by PostgreSQL, built authentication workflows, integrated payment-service logic, and maintained service and database behavior for NLP and speech products. Through teaching assistant roles at the University of Calgary, I mentored students in AWS-based data pipelines, cloud storage, large-scale data processing, SQL, debugging, and database coursework, which maps well to PLATO's focus on ETL development, pipeline reliability, data flow, and troubleshooting.

My machine learning and scientific-data background also helps me approach data quality carefully. I have built PyTorch workflows for digital elevation model super-resolution, evaluated models with terrain-aware metrics, and worked with validation issues where apparently good results could hide weak generalization. I know this role is ETL and data engineering rather than model research, but that experience helps me understand why clean data movement, validation, monitoring, and clear documentation matter.

I would bring a positive, self-managing, and communication-oriented approach to the team. My research and teaching work required clear written updates, careful documentation, and collaboration with technical and nontechnical audiences. I am comfortable learning new tools, asking precise questions, and taking ownership of small improvements while staying aligned with team standards and data governance expectations.

Thank you for considering my application. I would be grateful for the opportunity to discuss how my Python, backend, data pipeline, troubleshooting, and documentation experience can contribute to PLATO.

\closing{Sincerely,

Amir Mirzai Golpayegani

\href{mailto:amirgolpaa24@gmail.com}{amirgolpaa24@gmail.com}}

\end{letter}

\end{document}
